\documentclass[a4paper,11pt,oneside]{book}

\author{Vincenzo Maffione}
\title{Optimization in Virtual Machine Networking}
\date{\today}

%\usepackage{syntonly}
%\syntaxonly

\usepackage{lmodern} 

\usepackage{amssymb} % simboli matematici AMS
\usepackage{amsmath}

\usepackage{txfonts}

\hyphenation{pro-prie-t� in-di-vi-duo re-le-ga-re mo-del-lo mo-del-li o-pe-ra-zio-ni o-pe-ra-to-ri mo-du-lo ri-ve-la-to per-met-te-re rap-pre-sen-ta-zio-ne mo-del-liz-za-zio-ne ge-ne-ra-to o-biet-ti-vo li-ne-a-re no-no-stan-te po-si-ti-va e-la-bo-ra-zio-ne o-pe-ra-ti-vo sca-la-bi-li-t� i-ni-zia-liz-za-re in-di-vi-duo o-pe-ra-to-re va-lo-ri i-te-ra-zio-ni in-va-li-da-re pe-ri-fe-ri-che re-gi-stri ma-na-ge-ment e-si-ste pra-ti-ca va-ria-bi-li in-di-riz-zo o-pe-ra-zio-ne me-dian-te in-di-vi-duo sem-pli-ci-t� co-lo-ri ha-mil-to-nia-no que-sta in-clu-de-re ge-ne-tic}


\usepackage[english]{babel}

%\usepackage{indentfirst} % indentazione nella prima riga di ogni capoverso

\setlength{\parindent}{0pt} % elimina l'indentazione sulla prima riga di ciascun capoverso


%\usepackage[T1]{fontenc} % migliora la sillabazione di lingue non inglesi
\usepackage[latin1]{inputenc} %codifica per le lingue latine


% definizione dell'ambiente "abstract", per aggiungere il sommario
\usepackage{fancyhdr}
\newcommand{\fncyblank}{\fancyhf{}}
\newenvironment{abstract}%
{\cleardoublepage\fncyblank\null\vfill\begin{center}%
\bfseries\abstractname\end{center}}%
{\vfill\null}

\usepackage{graphicx}

\linespread{1.3}  % aumenta l'interlinea di un fattore 1.3


\begin{document}

%\maketitle

\begin{titlepage}
\pagestyle{empty}
\begingroup
% intestazione
\vspace*{-8.5\topskip}

\begin{figure}[!h]
\begin{center}
\includegraphics[scale=0.3]{logo_sssup}
\end{center}
\end{figure}

\vspace*{-0.5cm}

\begin{center}
        \large{\bf\expandafter\textsc {Universit\'a di Pisa}}\par 
		  \large{\bf\expandafter\textsc {Facolt\'a di Ingegneria}}\par
		  \large{\bf\expandafter\textsc {Laurea magistrale in Ingegneria Informatica}}\par
        \vspace{0.5cm}
        \large{\bf\expandafter\textsc {Tesi di Laurea}}
\end{center}



% Titolo
\vspace{1.5cm}
\begin{center}
        {\LARGE\bf\textit{Optimization in Virtual Machine Networking}}
\end{center}

\vspace{1.5cm}



\begin{center}
\begin{tabular}{l p{3.3cm} r c}

\textsc{Relatore} & & \textsc{Il candidato} \\
&&\\
\dotfill&&\dotfill \\
Prof. \textit{Luigi Rizzo} & &\textit{Vincenzo Maffione}\\
{\small Universit\'a di Pisa} & &\\
 & &  \\
\textsc{Relatore} & &  \\
 \dotfill&&\\
Prof. \textit{Giuseppe Lettieri} && \\
{\small Universit\'a di Pisa} && \\
 & & \\

\end{tabular}
\end{center}

\vspace{0cm}


\begin{center}
Anno Accademico 2012-2013
\end{center}
\par
\vfill\par 
\clearpage
\endgroup
  

\end{titlepage}

\newpage                                
\clearpage{\pagestyle{empty}\cleardoublepage}






\newpage\null\thispagestyle{empty}\newpage		%pagina bianca

\frontmatter					%Attiva numerazione romana perche' sei nell'intro


\begin{abstract}

Network performance is a critical aspect in Virtual Machine systems and its 
importance is becoming increasingly important in the world of computing.
These systems are commonly employed in the IT departments of several organizations, since they can help to 
build services that are highly reliabile, availabile and secure, improve efficiency
in computing resource usage, and so on.

\vspace{0.5cm}

In this thesis we are going to analize the state of the art of virtual machine networking, evaluating advantages
and drawbacks of the existing solutions.
We then propose a new approach, showing that with a small amount of code modifications, we can bring a classic emulated
network device (we take \texttt{e1000} as a reference example) to performances that are similar to the performances of 
paravirtualized solutions.

\vspace{0.5cm}

Lastly, we put together the lessons learned, and we introduce a new minimal paravirtualized solution, that can be implemented 
in total with about 2400 lines of code (driver + board emulation) and it is intended to overperform the currently existing solutions.

\end{abstract}


\newpage\null\thispagestyle{empty}\newpage		%pagina bianca


\tableofcontents
\listoffigures

\mainmatter					%Disattiva numerazione romana


\chapter{Introduction}

Standard computer systems are hierarchically organized as a three
layers stack. The lowest layer is the bare hardware, the middle one is the operating system and the 
upper layer contains the application software.
Two neighbour layers can communicate through a well-defined \emph{interface}, so that each layer can ignore how the
lower layers are actually implemented. In this way, the interface provides an \emph{abstraction} of the underlying software/hardware
resources.

This (relatively) simple architecture has proven to be very effective in delivering the IT services required
by companies, individual users and other organizations.

\vspace{0.5cm}

Neverthless, more complex computer systems organizations have been devised and used, in order to overcome some limitations
of the standard computer system model. These computing environments are are known as \emph{Virtual Machines} (VMs) systems.
By itself, the term \emph{Virtual Machine} can have several meanings, so when using it is important to point out what we are
addressing.
In general \emph{virtualization} provides a way of increasing the \emph{flexibility} of real hardware/software resources. When a 
resource  is virtualized, it appears to be a different virtual resource or even a \emph{set} of different virtual resources.

As an example, a single IA32 processor can be virtualized in a way that emulates more PowerPC processors. Once this is done, you
can build more standard 3-layer computer systems on top of each emulated PowerPC (virtual) processor, using unmodified OS and 
applications designed to be used on the PowerPC architecture.

In the VMs terminology, each virtualized system is called \emph{guest}, whereas the system providing the virtualization support
is called \emph{host}.
You can virtualize all the resources you wants: disks, network devices, memories and other peripherals.

\vspace{0.5cm}

Generally speaking VMs allow to build computer systems with more abstraction levels than the standard model have. This has important
advantages:
\begin{itemize}
  \item In terms of \emph{flexibility}, using VMs you can easily run programs compiled for a given Instruction Set Architecture (ISA) 
	and a given Operating System (OS) on top of a computer system that has a different ISA and/or a different OS. Using a standard
	system you would be bound to the ISA of your processor and the operating system installed on your machine.

  \item In terms of \emph{protection}, VMs can provide multiple isolated execution environments running on the same phisical machine.
	This allows to execute different applications in different VMs (each VM can have its own OS), so that if an application has a
	security hole, an attacker cannot use the hole to do malicious attacks to an applications running on a different VM. 
	This scenary is still possible when applications are run in the same OS.
  \item In terms of resources usage, VMs can help to reduce hardware costs and power consumption, since they 
\end{itemize}


\section{Virtual Machines}



%\begin{enumerate}
%	\item 24 \emph{tiles} (ossia mattonelle) che compongono il cluster: ogni tiles contiene due core IA
%	\item Una rete mesh composta da 24 router con picco di banda sul taglio pari a 256 GB/s
%	\item 4 DDR3 memory controllers integrati
%	\item supporto hardware per lo scambio dei messaggi
%\end{enumerate}

%\begin{figure}[bt]
%\centering
%\includegraphics[scale = 0.45]{SCCBlockDiagram.png}
%\caption{ Schema a blocchi dell'SCC }
%\label{fig:SCC}
%\end{figure}



\begin{thebibliography}{1}

\bibitem{ref:vmbook}
James E. Smith, Ravi Nair\\
Virtual Machines - Versatile Platforms for systems and processes\\
Elsevier, 2005\\

\bibitem{ref:x86-virt}
Keith Adams, Ole Agesen\\
A Comparison of Software and Hardware Techniques for x86 Virtualization\\
ACM, 2006\\

\bibitem{tassonomia}
Riccardo Poli, Mariusz Nowostawski (1999) \\
Parallel Genetic Algorithm Taxonomy

\bibitem{microprocessorReport}
Max Baron (2010) \\
The Single-chip Cloud Computer \\
Microprocessor report, www.MPRonline.com

\bibitem{intelDeveloperManual}
Intel (1995)
Pentium� Processor Family Developer's Manual, Architecture and Programming Manual

\bibitem{intelTechnicalReport}
Howard, Dighe et al. (2010) \\
A 48-Core IA-32 Message-Passing Processor with DVFS in 45nm CMOS

\bibitem{RCCEspecification}
Tim Mattson, Rob van der Wijngaart (2010) \\
RCCE: a Small Library for Many-Core Communication - Specification document

\bibitem{RCCEpaper}
Rob van der Wijngaart, Timothy Mattson, Werner Haas (2010) \\
Light-weight Communications on Intel's Single-Chip Cloud Computer Processor

\bibitem{SCCEAS}
Intel labs (2010) \\
SCC External Architecture Specification (EAS), Revision 1.1

\bibitem{SCCprogrammers}
Intel labs (2010) \\
The SCC Programmer's Guide, Revision 0.75

\bibitem{Baker}
Baker, James E. (1987)
Reducing Bias and Inefficiency in the Selection Algorithm
Proceedings of the Second International Conference on Genetic Algorithms and their Application (Hillsdale, New Jersey: L. Erlbaum Associates)

\end{thebibliography}


\end{document}

